% conclusions.tex (국문 v15, 2026-10-21): 목적, 근거, 의미의 순서.
\section{결론}\label{sec:conclusions}
ALT 관측은 드물고 군집되어 있으며, 이를 관측이 없는 셀과 해로 확장하는 모델은 계수 하나짜리 물리 모델의 전이성과 학습 지점에서의 기계학습 정확도 사이의 상충 관계에 놓인다. 이 연구는 학습기를 물리에 고정하면 그 상충 관계를 피할 수 있음을 보였다. 세 지역의 지점-연도 관측 \nncyAll{}개에서 각 관측 연도의 기후로 평가한 결과, 지역 라벨로 재보정한 Stefan 모델의 잔차를 학습한 gradient boosting 모델은 공간 블록 검증에서 Stefan 오차를 레나 델타에서 \ndrefLeblkRoabs{} cm, 캐나다에서 \ndrefCablkRoabs{} cm 줄였고, 판정되는 이득에 지역 라벨 25--100개만 필요하였으며, 자기 지역을 학습에서 뺐을 때도 Stefan 모델과 \ndrefLeregRoabs{} cm 안에 머물렀다. 같은 상황에서 직접 기계학습은 Stefan 오차를 \ndrefLeregDzabs{} cm 넘었다. 같은 지점과 해에서 채점한 공개 ML 제품은 자기 학습 지점 근처에서만 이 모델과 같은 정확도를 보였고 레나 델타의 얇은 활동층을 나타내지 못하였으며, 그곳에서 물리 잔차 모델이 제품 5종 모두보다 정확하였다. 실용적 결과는 세 가지이다. 적용 가능 영역과 예측 구간이 딸린 2011--2024년 연도별 1~km 지도, 모델과 제품의 정확도를 명시된 지점과 해에서 비교할 수 있게 하는 평가 프로토콜, 그리고 새 지역을 위한 규칙이다. 물리에서 출발하고, 지역 라벨이 계수 하나로는 지점을 잘못 나타낸다고 알려 주는 곳에 학습기를 더하며, 수천 개가 아니라 수십 개의 라벨에서 이득을 기대한다.
